%% TEMPORARY root file: compiles only Chapters 5 and 6 (with the thesis title page) to send for review.
%% The preamble is a copy of thesis.tex (10 October 2026). Delete this file and chapters_5_6.pdf when done.
%% Build from thesis/: pdflatex chapters_5_6 && bibtex chapters_5_6 && pdflatex chapters_5_6 && pdflatex chapters_5_6

%% (Master) Thesis template
% Template version used: v1.4
%
% Largely adapted from Adrian Nievergelt's template for the ADPS
% (lecture notes) project.


%% We use the memoir class because it offers a many easy to use features.
\documentclass[11pt,a4paper,oneside]{memoir}

%% Packages
%% ========

%% LaTeX Font encoding and language
\usepackage[T1]{fontenc} 
\usepackage[utf8]{inputenc}
\usepackage[english]{babel}

% Use Latin Modern (a standard, CM-like font family commonly used for
% academic papers). Computer Modern math will be used (good default math
% appearance for papers).
\usepackage{lmodern}
\usepackage{microtype}

%% [REC] TikZ is used to draw figures
\usepackage{tikz}
% For positioning library (helps position nodes more easily), 
% and arrows.meta for more arrow tips and styles
\usetikzlibrary{
  decorations.pathreplacing,
  calc,
  shapes,
  arrows.meta,
  positioning,
  backgrounds,
  fit,
  matrix,
  graphs,
  graphs.standard, 
  positioning, 
  arrows.meta,
  patterns,
  shapes.geometric,
  shapes.misc,
  decorations.pathmorphing,
  shadows,
}



%% The AMS-LaTeX extensions for mathematical typesetting.  Do not
%% remove.  
\usepackage{amsmath,amsfonts,mathrsfs}


%% NTheorem is a reimplementation of the AMS Theorem package. This
%% will allow us to typeset theorems like examples, proofs and
%% similar.  Do not remove.
%% NOTE: Must be loaded AFTER amsmath, or the \qed placement will
%% break
\usepackage[amsmath,thmmarks]{ntheorem}
 
%% LaTeX' own graphics handling
\usepackage{graphicx}

% Search for graphics in the `logos/` subdirectory as well
\graphicspath{{logos/}}

%% We unfortunately need this for the Rules chapter.  Remove it
%% afterwards; or at least NEVER use its underlining features.
\usepackage{soul} 

%% This allows you to add .pdf files. It is used to add the
%% declaration of originality.
\usepackage{pdfpages}

%% Some more packages that you may want to use.  Have a look at the
%% file, and consult the package docs for each.
%% See the TeXed file for more explanations

%% [OPT] Multi-rowed cells in tabulars
%\usepackage{multirow}


%% [OPT] Easily changeable quotes with \enquote{Text}
%\usepackage[german=swiss]{csquotes}

%% [REC] Format dates and time depending on locale
\let\ordinal\relax %surpresses a warning about ordinals
\usepackage{datetime}

%% [OPT] Provides a \cancel{} command to stroke through mathematics.
%\usepackage{cancel}

%% [NEED] This allows for additional typesetting tools in mathmode.
%% See its excellent documentation.
\usepackage{mathtools}

%% [ADV] Conditional commands
%\usepackage{ifthen}

%% [OPT] Manual large braces or other delimiters.
%\usepackage{bigdelim, bigstrut}

%% [NEED] Some extensions to tabulars and array environments.
\usepackage{array}

%% [OPT] Postscript support via pstricks graphics package. Very
%% diverse applications.
%\usepackage{pstricks,pst-all}

%% [?] This seems to allow us to define some additional counters.
%\usepackage{etex}

%% [ADV] XY-Pic to typeset some matrix-style graphics
%\usepackage[all]{xy}

%% [OPT] This is needed to generate an index at the end of the
%% document.
%\usepackage{makeidx}

%% [OPT] Fancy package for source code listings.  The template text
%% needs it for some LaTeX snippets; remove/adapt the \lstset when you
%% remove the template content.
\usepackage{listings}
\lstset{language=TeX,basicstyle={\normalfont\ttfamily}}

%% Python listings of Appendix E: colored syntax highlighting and line
%% numbers. xcolor is loaded by TikZ. Use with [style=python]; add float=htbp
%% to long listings so that they do not break across pages.
\definecolor{lstkeyword}{RGB}{0,51,179}
\definecolor{lstcvxpy}{RGB}{0,122,122}
\definecolor{lstcomment}{RGB}{0,128,0}
\definecolor{lststring}{RGB}{163,21,21}
\definecolor{lstnumber}{gray}{0.55}
\definecolor{lstframe}{gray}{0.75}
\definecolor{lstbackground}{gray}{0.97}
\lstdefinestyle{python}{
  language=Python,
  basicstyle=\ttfamily\small,
  columns=fullflexible,
  keepspaces=true,
  keywordstyle=\color{lstkeyword}\bfseries,
  emph={Variable,Parameter,Problem,Minimize,Maximize,sum,sum_squares,norm,solve,value},
  emphstyle=\color{lstcvxpy},
  commentstyle=\color{lstcomment}\itshape,
  stringstyle=\color{lststring},
  showstringspaces=false,
  numbers=left,
  numberstyle=\tiny\color{lstnumber},
  numbersep=10pt,
  xleftmargin=2em,
  frame=single,
  framesep=4pt,
  rulecolor=\color{lstframe},
  backgroundcolor=\color{lstbackground},
  breaklines=true,
  captionpos=b
}

%% [REC] Fancy character protrusion.  Must be loaded after all fonts.
\usepackage{microtype}

%% [REC] Nicer tables.  Read the excellent documentation.
\usepackage{booktabs}


%% [REC] Improves the spacing of \left \right
\usepackage{mleftright}

%% [REC] Improves spacing control.
\usepackage{xspace}

%% [REC] Bold math symbols
\usepackage{bm}

%% Our layout configuration.  DO NOT CHANGE.
\input{thesis_settings/layoutsetup.tex}

%% Theorem environments.  You will have to adapt this for a German
%% thesis.
\input{thesis_settings/theoremsetup.tex}

%% Helpful macros.
%% Custom commands
%% ===============

%% Special characters for number sets, e.g. real or complex numbers.
\newcommand{\C}{\ensuremath{\mathbb{C}}}
\newcommand{\K}{\ensuremath{\mathbb{K}}}
\newcommand{\N}{\ensuremath{\mathbb{N}}}
\newcommand{\Q}{\ensuremath{\mathbb{Q}}}
\newcommand{\R}{\ensuremath{\mathbb{R}}}
\newcommand{\Z}{\ensuremath{\mathbb{Z}}}
\newcommand{\X}{\ensuremath{\mathbb{X}}}

%% Fixed/scaling delimiter examples (see mathtools documentation)
\DeclarePairedDelimiter\abs{\lvert}{\rvert}
\DeclarePairedDelimiter\norm{\lVert}{\rVert}

% Define generic paired delimiters:
\DeclarePairedDelimiter{\Paren}{(}{)}
\DeclarePairedDelimiter{\Bracket}{[}{]}
\DeclarePairedDelimiter{\Brace}{\{}{\}}


% Big-O
\NewDocumentCommand{\Oh}{som}{%
  \ensuremath{\mathcal{O}\IfBooleanTF{#1}{%
    \Paren*{#3}%          % *-version => auto-size
  }{%
    \IfNoValueTF{#2}{%
      \Paren{#3}%         % no optional => normal size
    }{%
      \Paren[#2]{#3}%     % optional => \big, \Big, etc.
    }%
  }%
}}

%Big Omega:
\NewDocumentCommand{\Om}{som}{%
\ensuremath{\Omega\IfBooleanTF{#1}{%
    \Paren*{#3}%
  }{%
    \IfNoValueTF{#2}{%
      \Paren{#3}%
    }{%
      \Paren[#2]{#3}%
    }%
  }%
}}

% Big Theta:
\NewDocumentCommand{\Th}{som}{%
\ensuremath{\Theta\IfBooleanTF{#1}{%
    \Paren*{#3}%
  }{%
    \IfNoValueTF{#2}{%
      \Paren{#3}%
    }{%
      \Paren[#2]{#3}%
    }%
  }%
}}

% little-o:
\NewDocumentCommand{\oh}{som}{%
\ensuremath{o\IfBooleanTF{#1}{%
    \Paren*{#3}%
  }{%
    \IfNoValueTF{#2}{%
      \Paren{#3}%
    }{%
      \Paren[#2]{#3}%
    }%
  }%
}}

% little-omega:
\NewDocumentCommand{\om}{som}{%
\ensuremath{\omega\IfBooleanTF{#1}{%
    \Paren*{#3}%
  }{%
    \IfNoValueTF{#2}{%
      \Paren{#3}%
    }{%
      \Paren[#2]{#3}%
    }%
  }%
}}

% Probability:  \Prob{Event} =>  P{Event},   \Prob*{...} => auto-sized { ... }
%    By default, we use curly braces for probabilities.
\NewDocumentCommand{\Prob}{som}{%
\ensuremath{\mathbb{P}% or \Pr, if you prefer
  \IfBooleanTF{#1}{%
    % *-version => auto-sized braces
    \Bracket*{#3}%
  }{%
    % no star
    \IfNoValueTF{#2}{%
      % no optional argument => normal size braces
      \Bracket{#3}%
    }{%
      % optional argument => \big, \Big, etc.
      \Bracket[#2]{#3}%
    }%
  }%
}}

% Expectation:  \E{X} => E[X],   \E*{X} => auto-sized [ X ]
\NewDocumentCommand{\E}{som}{%
\ensuremath{\mathbb{E}
  \IfBooleanTF{#1}{%
    \Bracket*{#3}%
  }{%
    \IfNoValueTF{#2}{%
      \Bracket{#3}%
    }{%
      \Bracket[#2]{#3}%
    }%
  }%
}}

\NewDocumentCommand{\Max}{s o m m}{%
\ensuremath{\max
  \IfBooleanTF{#1}{%
    \Brace*{#3,\, #4}%
  }{%
    \IfNoValueTF{#2}{%
      \Brace{#3,\, #4}%
    }{%
      \Brace[#2]{#3,\, #4}
    }%
  }%
}}

\NewDocumentCommand{\Min}{s o m m}{%
\ensuremath{\min
  \IfBooleanTF{#1}{%
    \Brace*{#3,\, #4}%
  }{%
    \IfNoValueTF{#2}{%
      \Brace{#3,\, #4}%
    }{%
      \Brace[#2]{#3,\, #4}
    }%
  }%
}}

% Exponential:  \Exp{x} => exp(x)
\NewDocumentCommand{\Exp}{som}{%
\ensuremath{\exp
  \IfBooleanTF{#1}{%
    \Paren*{#3}%
  }{%
    \IfNoValueTF{#2}{%
      \Paren{#3}%
    }{%
      \Paren[#2]{#3}%
    }%
  }%
}}

% You might want \Log, \Var, \Cov, etc. Similarly:
\NewDocumentCommand{\Log}{som}{%
\ensuremath{\log
  \IfBooleanTF{#1}{%
    \Paren*{#3}%
  }{%
    \IfNoValueTF{#2}{%
      \Paren{#3}%
    }{%
      \Paren[#2]{#3}%
    }%
  }%
}}

\NewDocumentCommand{\Var}{som}{%
\ensuremath{\mathrm{Var}
  \IfBooleanTF{#1}{%
    \Bracket*{#3}%
  }{%
    \IfNoValueTF{#2}{%
      \Bracket{#3}%
    }{%
      \Bracket[#2]{#3}%
    }%
  }%
}}


\NewDocumentCommand{\SetBuilder}{s o m m}{%
\ensuremath{\IfBooleanTF{#1}{%
    % Star version: auto-sized braces
    \Brace*{\,#3 : #4\,}%
  }{%
    % Non-star version
    \IfNoValueTF{#2}{%
      % No optional argument => normal-size braces
      \Brace{\,#3 : #4\,}%
    }{%
      % Optional argument => user-specified size, e.g. \big, \Big, \bigg
      \Brace[#2]{\,#3 : #4\,}%
    }%
  }%
}}



%% The xspace package is used to automatically handle spacing, regardless of whether the command is followed by punctuation or not.
\newcommand{\onemax}{\textsc{OneMax}\xspace}


%% \eps is quicker to type and the \varepsilon font is more common and nice.
\newcommand{\eps}{\ensuremath\varepsilon}

%% Also set the alternate phi as default, no real reason to use the old one, but one can.
\let\oldphi\phi
\renewcommand{\phi}{\ensuremath{\varphi}}

%% --- Notation shortcuts (from user) ---
\newcommand{\1}{\bm{1}}
\newcommand{\T}{\mathsf{T}}
\newcommand{\had}{\odot} % Hadamard product


%% --- Math Operators ---
\DeclareMathOperator*{\argmin}{arg\,min}
\DeclareMathOperator*{\argmax}{arg\,max}


%% Make document internal hyperlinks wherever possible. (TOC, references)
\usepackage[linkcolor=black,colorlinks=true,citecolor=black,filecolor=black]{hyperref}

% Load cleveref after hyperref
\usepackage{cleveref}

\usepackage{upgreek}

\usepackage{natbib}
\usepackage{multirow}

% We now name most things that are to be referred to.
% You can add similar commands for your own defined environments
% THEOREM
\crefname{theorem}{theorem}{theorems}
\Crefname{theorem}{Theorem}{Theorems}

% LEMMA
\crefname{lemma}{lemma}{lemmas} 
\Crefname{lemma}{Lemma}{Lemmas}

% COROLLARY
\crefname{corollary}{corollary}{corollaries}
\Crefname{corollary}{Corollary}{Corollaries}

% PROPOSITION
\crefname{proposition}{proposition}{propositions}
\Crefname{proposition}{Proposition}{Propositions}

% CLAIM
\crefname{claim}{claim}{claims}
\Crefname{claim}{Claim}{Claims}

% DEFINITION
\crefname{definition}{definition}{definitions}
\Crefname{definition}{Definition}{Definitions}
% EXAMPLE
\crefname{example}{example}{examples}
\Crefname{example}{Example}{Examples}
%% The AMS-LaTeX extensions for mathematical typesetting.  Do not
%% remove. We avoid loading amssymb directly because `newtxmath` will
%% provide compatible symbol definitions; this prevents a \Bbbk clash.
\usepackage{amsmath,amsfonts,mathrsfs}
% REMARK
\crefname{remark}{remark}{remarks} 
\Crefname{remark}{Remark}{Remarks}

% RULE - This one is special
\crefname{Rule}{rule}{rules}
\Crefname{Rule}{Rule}{Rules}

 
%% [REC] This is for restating theorems.
\usepackage{thmtools}
\usepackage{thm-restate}      


%% [REC] packages for pseudocode
\usepackage{algorithm}        % A float wrapper for algorithms
\usepackage{algorithmicx}
\usepackage{algpseudocode}


%% Document information
\title{Multi-Currency Portfolio Optimization with General Risk
Measures} 
\author{Anej Rozman}
\thesistype{Master Thesis} 
\advisors{Supervisor: Prof.\ Patrick Cheridito \\ Co-Supervisor: Dr.\ Urban Ulrych} 
\department{Department of Mathematics}   
\date{} 

%% Settings for theorem D.3 in appendix and chapter 6.2
\usepackage{cleveref}

\theoremstyle{nonumberplain}
\theoremheaderfont{\bfseries}
\theorembodyfont{\itshape}
\theoremsymbol{}
\newtheorem{demeaningthmDthree}{Theorem D.3}
\theoremstyle{plain}
\theoremsymbol{}

%% Cross-references to parts of the thesis that are not included in this file.
%% Copied from thesis.aux (full build of 10 October 2026), so they show the thesis numbering.
\newlabel{app:demeaning}{{D}{59}{Coherent Risk Measures and Deviation Measures}{appendix.D}{}}
\newlabel{app:demeaning@cref}{{[chapter][4][]D}{[1][59][]59}{}{}{}}
\newlabel{chap:data-and-implementation}{{4}{7}{Data and Implementation}{chapter.4}{}}
\newlabel{chap:data-and-implementation@cref}{{[chapter][4][]4}{[1][7][]7}{}{}{}}
\newlabel{chap:optimization-framework}{{3}{3}{Optimization Framework}{chapter.3}{}}
\newlabel{chap:optimization-framework@cref}{{[chapter][3][]3}{[1][3][]3}{}{}{}}
\newlabel{sec:backtest-configuration}{{4.3}{12}{Backtest Configuration}{section.4.3}{}}
\newlabel{sec:backtest-configuration@cref}{{[section][3][4]4.3}{[1][12][]12}{}{}{}}
\newlabel{sec:mean-risk-portfolio-optimization}{{3.2}{6}{Mean--Risk Portfolio Optimization}{section.3.2}{}}
\newlabel{sec:mean-risk-portfolio-optimization@cref}{{[section][2][3]3.2}{[1][5][]6}{}{}{}}
\newlabel{thm:demeaning}{{D.{3}}{60}{Coherent Risk Measures and Deviation Measures}{theorem.D.3}{}}
\newlabel{thm:demeaning@cref}{{[theorem][3][4]D.{3}}{[1][59][]60}{}{}{}}

%% ======================================================= %%
\begin{document}

\frontmatter

%% Title page (same as thesis.tex)
\begin{titlingpage}
  \calccentering{\unitlength}
  \begin{adjustwidth*}{\unitlength-24pt}{-\unitlength-24pt}
    \maketitle
  \end{adjustwidth*}
\end{titlingpage}

\mainmatter

%% Keep the thesis numbering, so the first chapter here is Chapter 5
\setcounter{chapter}{4}
\chapter{Backtest Analysis}\label{chap:backtest-analysis}



\section{Baseline Backtests} \label{sec:baseline-backtests}

WHAT TO COVER:

Define the baselin

three metric tables, one per universe size

three cumulative return plots, one per universe size

First, read tables. which measure is the best performer, where does joint beat overlay and by how much, where is it the reverse

Second, the universe size pattern. At top-10 the joint and overlay results differ most, and the direction of the difference depends on the measure. At top-100 they converge.

Third, the main claim. On realized Sharpe, joint and overlay are close across configurations, and no single risk measure delivers a consistent advantage for joint. This sets up what the modifications in the next two sections do or do not change.

\section{The Effects of Regularization} \label{sec:effects-of-regularization}

Opens with a one-paragraph lead stating that this section repeats the baseline comparison with the regularized versions of the joint and overlay programs. Point back to the Optimization Framework chapter for the regularization scheme.

Rather than repeating the three full tables, show one comparison figure and one summary table. The figure has three panels, one per universe. Each panel shows regularized vs unregularized Sharpe (or your headline metric) for joint and overlay across the five measures. Simple bar chart or dot plot with lines connecting the regularized and unregularized points, so the reader sees the direction and size of the change in one glance.

The summary table shows the delta: regularized minus unregularized, by strategy and measure, at each universe. One row per measure, columns grouped by strategy, three sub-columns per group (one per universe). Negative deltas mean regularization hurt Sharpe, positive means it helped.

The full regularized tables with all metrics go in the appendix.

What to write. Two or three paragraphs.

First, what regularization does to the levels. Does it improve Sharpe, hurt it, or move it only slightly. State the direction cleanly.

Second, what it does to the joint-overlay gap. If regularization narrows the gap, say so and give the numbers. If it leaves the gap unchanged, say that. This is the main point of the section because it foreshadows the next chapter's mechanism argument.

Third, any measure-specific or universe-specific surprises. If regularization helps the tail measures more than the deviation measures, or vice versa, say so. If one universe behaves differently, flag it. Do not force a story that is not there.

\section{The Effects of Adding a Mean Term} \label{sec:effects-of-adding-a-mean-term}

Opens with a one-paragraph lead stating that this section adds a mean term to the unregularized joint and overlay programs, repeating the baseline comparison. Point to the Optimization Framework chapter for how the mean term is specified and calibrated.

Same visual structure as Section 2: one comparison figure (three panels, one per universe) showing with-mean-term vs no-mean-term Sharpe for joint and overlay across measures, plus a delta table.

Then the extra analysis you have in mind. A second table that groups the delta by deviation measures vs tail measures. Average delta Sharpe for the three deviation measures and for the two tail measures, by strategy and universe. This is the table that shows the group split. If the mean term helps the deviation measures more than the tail measures (or less), it will be visible here.

The full tables go in the appendix.

What to write. Three paragraphs.

First, levels. Does adding the mean term improve realized performance overall. State the direction. The expectation is that it helps, since the mean term adds a return-seeking component to a pure risk minimizer. If it does not, say so and discuss why briefly.

Second, the deviation vs tail split. The tail measures are already translation-equivariant and carry an implicit mean term, so adding an explicit mean term with your calibrated kappa should strengthen an existing channel. The deviation measures do not have this channel, so the mean term is giving them something new. Predict what you expect to see (deviation measures benefit more if the new channel is strong, less if the existing tail-measure mean channel was already most of what mattered), then say which the data support.

Third, joint vs overlay with the mean term. Does the mean term widen or narrow the joint-overlay gap. The key claim here is that if the mean term helps joint and overlay by similar amounts, the "joint lets you exploit covariance with the mean" story is weak, because that story predicts joint benefits more. This matters because it closes off the obvious defense of joint optimization. State the result cleanly.

\section{Conclusion} \label{sec:backtest-conclusion}


\section{Backtest Results of Expected Shortfall Optimization Approaches}

We can see in table \ref{tab:es-universe-sizes} that the performance of the expected shortfall optimization approaches varies across 
different universe sizes. In the small universe, the optimized portfolios generally outperform the equal weight benchmark in terms of 
risk-adjusted returns, as measured by the Sharpe Ratio and Sortino Ratio. However, as we increase the universe size to medium and large, 
the performance of the optimized portfolios relative to the equal weight benchmark deteriorates, with some strategies even underperforming 
the benchmark in terms of risk-adjusted returns. This suggests that the benefits of optimization may diminish as the investment universe expands, 
potentially due to increased estimation error and model uncertainty in larger universes. 

\newpage
% checklist of what to compare 
List of results interpretations to write up
\begin{itemize}
  \item \textbf{Optimizers vs 1/N} All strategies outperform 1/N on all risk metrics and SR, SOR and ESR in all 3 universe sizes. The risk reduction is larger in ES and Vol than in MD, suggesting that the optimized portfolios are better at controlling tail risk and volatility, but may still be sucseptible to large drawdowns.
  
  \item \textbf{Risk metrics with increasing universe size} 
  We can observe that the risk metrics for the 1/N portfolio do not improve by much as the universe size increases. They do increase substantially for the optimized portfolios. They drop by around 50\% from the small to the large universe, while the 1/N portfolio only drops by around 15\%. This suggests that the benefits of optimization are more pronounced in larger universes, where there is more opportunity to diversify and reduce risk.

  \item \textbf{Joint-Reg vs.\ Overlay-Reg} In the small universe, Joint-Reg is ahead on SR, SOR and ESR  at slightly higher. Still waiting for results on the medium and large universes. Hypothesis: with the large universe there is no more benefit to the joint optimization. 

  \item \textbf{Regularized vs.\ unregularized.} Overlay vs.\ Overlay-Reg and Joint vs.\ Joint-Reg. Regularization does not improve risk metrics but does improve ER, SOR and ESR. It typically reduces turnover, but not in every situation. 

  \item \textbf{Plots} TODO: Compare drawdown depth of the optimizers against 1/N around 2002, 2008--09 and 2020, and compare how that gap changes based on universe size.

\end{itemize}

\newpage


\section{Backtest Results of Mean-Variance Optimization Approaches}

\section{Backtest Results of Mean Absolute Deviation Optimization Approaches}

FINISH WITH THIS SO YOU MOTIVATE THE NEXT CHAPTER:

The backtests of \Cref{chap:backtest-analysis} showed no consistent return advantage of
joint optimization over the sequential currency overlay. Across the five risk measures, the three universe sizes
and both program families, neither approach reliably outperforms the other in terms of raw returns. Both programs are risk minimizers, so the question that remains is whether solving for asset and
currency weights simultaneously delivers a lower value of the risk measure being minimized, and whether that advantage is consistenly realized out-of-sample. 
\chapter{Joint versus Sequential Optimization}\label{chap:joint-vs-sequential}

% #############################################################################
% WRITING GUIDE (delete before submission)
%
% THE CHAPTER'S CLAIM -- everything in here serves this one sentence:
%
%   The joint program's in-sample advantage over the sequential overlay is
%   inflated by estimation error in proportion to how little of the sample the
%   risk measure uses. The in-sample gap therefore measures overfitting
%   capacity rather than economic value, and the measures with the strongest
%   in-sample case deliver the least out of sample.
%
% Note the order of that sentence. The MECHANISM is the claim; the inversion
% across risk measures is its consequence. This matters for how you defend the
% chapter: the mechanism is supported by the level sweep (14 levels, one
% measure, slopes with confidence intervals) and by the universe sweep, both of
% which are direct evidence. The cross-measure inversion rests on five points
% and is corroboration. Do not lead with the five-point result and do not let
% the chapter depend on it.
%
% PAGE BUDGET (~16 pages)
%   6.1  In-sample case ................ 2 pages
%   6.2  What survives ................. 2-3 pages
%   6.3  Mechanism ..................... 6 pages   <- the chapter
%   6.4  The inversion ................. 2 pages
%   6.5  What does not repair it ....... 2 pages
%   6.6  When justified ................ 1-2 pages
%
% Note 6.3 and 6.4 are SWAPPED relative to the first draft of this outline:
% mechanism first, inversion as its consequence.
%
% TWO THINGS TO FIX BEFORE WRITING
%
% (a) SIGN CONVENTION. Your draft in the empirical chapter defines
%     g = joint - overlay, so NEGATIVE means joint is better. The code, every
%     figure and every table define gap = overlay - joint, so POSITIVE means
%     joint is better. Adopt the code's convention and fix the prose.
%
% (b) MATERIAL TO MIGRATE from the end of
%     chapters/empirical_analysis_and_baseline_results.tex:
%       - the hypothesis paragraphs  -> Section 6.2
%       - the currency-share arm     -> Section 6.2
%       - the decomposition derivation g^oos = g^is + Delta_epsilon
%                                    -> Section 6.1, as the formal backbone
%     That derivation is good and does real work. It already says what this
%     chapter proves empirically.
% #############################################################################

% TODO: roadmap paragraph, in the style of the other chapter openers. One
% paragraph, naming each section and what it establishes, in order.


\section{The In-Sample Case for Joint Optimization}\label{sec:joint-in-sample-case}

% ---------------------------------------------------------------------------
% THE CLAIM
%   Joint optimization weakly dominates the sequential overlay in sample by
%   construction, and the magnitude of that dominance differs systematically
%   across risk measures -- it is largest for the tail measures.
%
% WHY IT IS TRUE
%   1. The nesting argument, stated as a proposition in Chapter 3 and recalled
%      here. The overlay's solution is feasible for the joint program, and the
%      overlay's second step is a partial minimisation of the same objective:
%
%        f(theta_J) = min_{x,psi} f(x,psi) <= min_psi f(x_O, psi) = f(theta_O)
%
%      Use the derivation already drafted in the empirical chapter.
%   2. Then the decomposition, which is the chapter in one equation:
%
%        g^oos = [f(theta_J) - f(theta_O)] + [eps(theta_J) - eps(theta_O)]
%                \_______ g^is _______/     \____ Delta_eps >= 0 ____/
%
%      where eps(theta) = f^oos(theta) - f(theta) is the generalization gap.
%      The second term is non-negative in expectation because theta_J is
%      selected over a strictly larger set and therefore fits more estimation
%      noise. State this as the chapter's organising identity: the same
%      freedom that wins in sample is what costs out of sample. Everything
%      after this section measures Delta_eps.
%   3. Empirical verification: gap_in < 0 on 8 of 1365 unregularized
%      date-cells (0.6%), all at magnitude ~1e-7. Say explicitly that this is
%      solver tolerance, not violation, and that it doubles as an
%      implementation correctness check.
%   4. The in-sample gap by measure, averaged over the three universes:
%
%        SOES 7.23%   ES 3.71%   SMV 2.13%   MAD 1.91%   MV 1.45%
%
% WHAT THIS SECTION MUST SET UP
%   The expectation that the tail measures are the most promising place to
%   apply joint optimization. State that expectation plainly -- it is the one
%   Section 6.4 overturns, and it has to be on the page to be overturned.
%
% DATA: results/in_sample_vs_out_of_sample/gap_summary.csv (unregularized),
%       gaps.csv for the violation count.
% ---------------------------------------------------------------------------


\section{What Survives Out of Sample}\label{sec:what-survives}

% ---------------------------------------------------------------------------
% THE CLAIM
%   The in-sample advantage does not survive out of sample except in the
%   smallest universe, and it declines monotonically in the number of assets.
%
% WHY IT IS TRUE
%   Out-of-sample relative gap with bootstrap confidence intervals. Significant
%   at top-10 for ES, MAD, MV and SMV; at top-50 only for MV and SMV; at
%   top-100 for nothing; and for SOES nowhere at all.
%
% THE HYPOTHESIS (from the migrated text)
%   The joint program must estimate the asset-currency dependence structure
%   directly; the overlay sidesteps it by sequencing, which is a form of
%   dimensionality reduction. So joint should win where estimation is easy and
%   lose where it is not.
%
% THE OBJECTION THIS SECTION MUST SURVIVE
%   Is the decline driven by the currency share or by the asset count? They
%   move together in the main grid. The K-reduction arm separates them: hold
%   the asset universe fixed and cut the tradable currency block to K = 10, 6
%   (top-50) and K = 10, 6, 3 (top-10). Because s = K/(K + 14N), top-10 at
%   K = 3 lands at s = 0.021 against top-50's 0.020 with 140 securities rather
%   than 700. Keep this to a paragraph and a figure -- it is a robustness check
%   on the hypothesis, not a finding.
%
% FIGURE: results/in_sample_vs_out_of_sample/gap_vs_currency_share.pdf
%   (currently fig:gap-vs-currency-share in the empirical chapter; move the
%   float here rather than duplicating it).
%
% DO NOT interpret the cross-measure differences yet. This section establishes
% that the advantage decays; Section 6.3 explains why; Section 6.4 shows the
% decay is systematic across measures. Keep them apart or the argument blurs.
% ---------------------------------------------------------------------------


\section{Mechanism: The Information Budget}\label{sec:mechanism}

% ---------------------------------------------------------------------------
% THIS IS THE CHAPTER. Six pages. Three converging tests of one mechanism.
%
% THE CLAIM
%   Delta_eps -- the excess generalization gap the joint program pays for its
%   larger feasible set -- grows as the number of observations the risk measure
%   uses falls. The in-sample advantage is inflated by exactly that amount.
%
% Open the section by stating the mechanism in terms of an information budget:
% the joint program estimates N + K parameters; what varies across the three
% tests below is how many observations the objective actually uses to do it.
% Each test varies one thing and holds the others fixed. Say that explicitly
% before the subsections -- it is what makes three results into one argument.
% ---------------------------------------------------------------------------

\subsection{Across Tail Depth Within One Risk Measure}\label{sec:mechanism-level}

% ---------------------------------------------------------------------------
% PUT THIS FIRST. It is the cleanest test and the chapter's strongest evidence:
% it holds the risk measure, the universe, the constraint set and the data
% fixed, and varies only how much of the sample the objective uses.
%
% THE CLAIM
%   Within a single risk measure, looking deeper into the tail raises the
%   in-sample advantage without raising the realized one. The additional
%   in-sample advantage is entirely estimation error.
%
% THE DESIGN
%   ES swept over 14 confidence levels (50, 60, 70, then 2.5-point steps from
%   72.5 to 97.5) on all three universes, 42 cells. Tail observations in the
%   estimation window fall from 1044 at level 60 to 65 at level 97.5 while the
%   number of free weights is unchanged. That ratio IS the information budget,
%   and the sweep varies it by a factor of 16 with nothing else moving.
%
% THE NUMBERS
%   Slope of the relative gap per percentage point of confidence level:
%
%                   in-sample                    out-of-sample (annual)
%     top-10    +0.00162 [0.00135, 0.00186]   +0.00005 [-0.00063, 0.00076]
%     top-50    +0.00100 [0.00078, 0.00124]   +0.00016 [-0.00029, 0.00064]
%     top-100   +0.00060 [0.00038, 0.00087]   +0.00037 [-0.00006, 0.00079]
%
%   In-sample: significantly positive at every universe. Out-of-sample: zero
%   at every universe. At top-10 the in-sample gap runs from 3.11% at level 60
%   to 10.59% at 97.5 while the realized gap sits at 2.2-2.8% throughout.
%
%   The overfit ratio (realized decay relative to 1/N's, which optimizes
%   nothing and so isolates estimation error) rises in the same shape:
%   top-10 from 1.34 at level 60 to 2.01 at 97.5; top-100 from 3.07 to 5.88.
%
% THE SENTENCE TO WRITE
%   The in-sample gap decomposes into a level-independent structural component
%   (~2.5% at top-10, which survives) and a level-increasing estimation
%   component (which does not). Deeper tail targeting buys no additional real
%   advantage; it buys the joint program more room to overfit a thinning tail.
%
% THE OBJECTION THIS MUST SURVIVE -- and it is a real one, so do not bury it.
%   f_out is scored on the quarter actually held, whose tail thins to ~3
%   observations at level 95 and ~2 at 97.5. At that point f_out is an extreme
%   order statistic, not an expectation. A second horizon (the following year,
%   ~252 days) is therefore scored alongside it, lifting the level-97.5 tail
%   to ~6 observations. The quarterly horizon suggests the realized gap
%   DECLINES at high levels; the annual horizon shows it is FLAT and agrees
%   with the pooled production backtests where the quarterly one does not.
%   Report both, say which you rely on and why. Handling this openly is what
%   makes the section credible; an examiner will go looking for it.
%
% DATA: results/in_sample_vs_out_of_sample_es_at_diff_levels/
% ---------------------------------------------------------------------------

\subsection{Across Problem Dimension}\label{sec:mechanism-dimension}

% ---------------------------------------------------------------------------
% THE CLAIM
%   Holding the risk measure fixed and raising the number of assets has the
%   same effect: more parameters against the same observations, so more of the
%   in-sample advantage is estimation error.
%
% THE NUMBERS
%   Overfit ratios for ES rise from 1.38 at top-10 to 3.09 at top-100. Both
%   programs overfit more as the universe grows, and the joint program
%   overfits more than the overlay in 12 of 15 cells (paired, CIs clear of
%   zero). The three exceptions -- top-10 at levels 60, 70, 80 -- are exactly
%   the cells where the joint advantage survives out of sample.
%
% THE SENTENCE TO WRITE
%   The joint program's advantage survives if and only if it does not overfit
%   more than the overlay does. That is a mechanism, not a correlation, and it
%   is worth stating as the section's punchline.
%
% Note that both programs still beat 1/N in absolute realized risk everywhere
% (0.38 to 0.69 of it). Say so -- it prevents the reader concluding that
% optimization does not work, when the claim is narrower.
% ---------------------------------------------------------------------------

\subsection{Across Risk Measures}\label{sec:mechanism-measures}

% ---------------------------------------------------------------------------
% THE CLAIM
%   The same mechanism operates across risk measures. Deviation measures
%   (variance, MAD, semi-variance) use every observation; tail measures (ES,
%   SOES) use only the tail. The tail measures therefore have the largest
%   in-sample gap and the least realizable one.
%
% THE COMPETING EXPLANATION THIS MUST RULE OUT
%   ES and SOES are not location-free -- they carry an implicit mean. Is their
%   behaviour driven by that rather than by tail focus? Test by de-meaning,
%   which makes them location-free. Result: the in-sample ceiling falls but
%   does not vanish. So the implicit mean explains part and the tail focus the
%   rest. Report the magnitude of the fall; this is the experiment that
%   converts a plausible story into a tested one.
%
% DATA: results/in_sample_vs_out_of_sample_demeaned/
% ---------------------------------------------------------------------------


\section{The Inversion Across Risk Measures}\label{sec:the-inversion}

% ---------------------------------------------------------------------------
% THE CLAIM
%   Because the in-sample gap is inflated in proportion to the information
%   deficit, its magnitude predicts the realized benefit NEGATIVELY across
%   risk measures.
%
% This section is now a CONSEQUENCE of Section 6.3, not an independent
% finding. Write it that way: "the mechanism implies an ordering, and the
% ordering is observed." That framing is both more honest and more persuasive
% than presenting it as a standalone discovery.
%
% THE NUMBERS -- averaged over the three universes, unregularized:
%
%     measure    in-sample   out-of-sample   share realized
%     SOES         7.23%        0.47%             7%
%     ES           3.71%        1.14%            31%
%     SMV          2.13%        1.87%            88%
%     MAD          1.91%        1.18%            62%
%     MV           1.45%        1.52%           105%
%
%   Spearman rank correlation between the in-sample and out-of-sample gap
%   across the five measures: -0.50 (top-10), -0.20 (top-50), -0.60
%   (top-100). Negative at every universe size.
%
% THE HONESTY REQUIREMENT -- do not skip this, it is what makes the claim
% survive a viva:
%   Five measures is a thin basis for a rank correlation and the correlation
%   is not itself a powered test. Say so in one sentence, and point to
%   Section 6.3.1: the level sweep reproduces the same effect WITHIN one
%   measure with 14 points and confidence intervals, which is what makes the
%   cross-measure ordering credible rather than a coincidence over five
%   points. Presented in that order the thin result is corroboration of a
%   well-supported mechanism instead of a load-bearing claim.
%
% FIGURE: not yet built. Whatever form you choose, it has to put the
% in-sample magnitude and the realized share in the same coordinate system.
% The existing five-panel figure cannot express this claim because neither
% variable of the claim is on an axis.
% ---------------------------------------------------------------------------


\section{What Does Not Repair It}\label{sec:what-does-not-repair}

% ---------------------------------------------------------------------------
% Defence. Pre-empts the two obvious objections. Two pages, no more.
% ---------------------------------------------------------------------------

\subsection{Regularization}\label{sec:no-repair-regularization}

% ---------------------------------------------------------------------------
% METHODOLOGICAL POINT FIRST -- it justifies the scope of the whole chapter.
%   The in-sample gap framework does NOT extend to penalised programs. The
%   nesting proposition requires the scorer to be the objective. Regularized
%   joint and overlay programs minimise Risk + penalty with penalties selected
%   per date that differ between the two, so they optimise different
%   objectives while being scored on a third. Nothing is bounded and the
%   in-sample/out-of-sample difference is not the decay of anything.
%
%   EVIDENCE THAT THIS IS NOT HYPOTHETICAL: the ceiling gap_in >= 0 is
%   violated on 30.0% of regularized date-cells (MV 44.7%, SMV 38.5%,
%   MAD 20.1%, ES 16.8%) against 0.6% at solver tolerance for unregularized.
%   Cite these numbers. They turn the restriction of this chapter to the
%   unregularized family from an apparent omission into a consequence of the
%   proposition -- which is a much better thing for an examiner to find.
%
% THE CLAIM, answered out of sample only, by comparing realized risk of the
% production backtests (no in-sample quantity, no ceiling required):
%   Regularization reduces the absolute estimation error both programs suffer
%   without changing which of them is lower.
%
%   ES at level 80, relative reduction in realized ES_80, joint vs overlay:
%
%     universe      unregularized     regularized
%     top-10          +3.00% (sig)     +2.01% (sig)
%     top-50          +1.76% (sig)     +1.40% (sig)
%     top-100         +0.15% (ns)      -0.21% (ns)
%
%   Same pattern, same cut-off, slightly smaller under regularization. And the
%   overfit ratio falls sharply: at top-100 from 3.09 to 1.93, with realized
%   risk relative to 1/N improving from 0.384 to 0.350. Regularization pulls
%   both programs toward the same floor without changing their ordering.
% ---------------------------------------------------------------------------

\subsection{A Mean Term in the Objective}\label{sec:no-repair-mean}

% ---------------------------------------------------------------------------
% THE CLAIM
%   Adding a mean term helps the deviation measures substantially and does
%   little for ES and SOES, which already carry the mean implicitly.
%
% Frame as consistent with the mechanism: the measures with information to
% spare can afford to estimate a mean; the information-starved ones cannot.
%
% The nesting proposition fails here too, and for the same reason -- gamma is
% selected per date and the joint and overlay programs routinely select
% different values, so the two minimise different objectives. Say so and rely
% on realized quantities.
%
% DATA: results/in_sample_vs_out_of_sample_w_mean/, top_runs_w_mean.csv
% ---------------------------------------------------------------------------


\section{When Joint Optimization Is Justified}\label{sec:when-justified}

% ---------------------------------------------------------------------------
% THE CLAIM
%   Joint optimization is justified in small universes with variance-based
%   measures, and specifically not with the tail measures that look most
%   promising in sample.
%
% State that this is counterintuitive and that it follows directly from the
% mechanism: the tail measures' in-sample appeal IS their estimation problem.
%
% MAGNITUDE, net of transaction costs (this matters -- the gap analysis
% ignores costs and the production comparison does not, which closes the
% objection that higher currency turnover eats the benefit):
%   At top-10, joint reduces realized volatility by 2.5-5.7% and realized ES
%   by 1.8-3.9% relative to the overlay, with no detectable effect on mean
%   return. At top-100 nothing is significant on any measure.
%
% THE SHARPE QUESTION, and why it is the wrong instrument:
%   Sharpe divides an unestimable numerator by the improved denominator. The
%   implied Sharpe change from a 2.5% volatility reduction is ~0.003 against a
%   detection threshold of ~0.15 at this sample length. Report it as a bound
%   ("we can exclude a Sharpe improvement above roughly 0.10") rather than as
%   a null. Note that the same bootstrap on the same return series DOES detect
%   the volatility and ES effects, so the null is not a dead instrument -- it
%   is a statement about what Sharpe can resolve.
%
% THE HONEST LIMIT -- state it plainly, hedging would weaken the chapter:
%   At the universe sizes relevant for international equity portfolios
%   (700-1400 securities against 14 currencies), the two approaches are
%   statistically indistinguishable out of sample, and the sequential overlay
%   is cheaper to compute. That is the practical conclusion.
% ---------------------------------------------------------------------------


\backmatter

\bibliographystyle{plain}
\bibliography{refs}

\end{document}
